\section{Deterministic LP formulation}
The decision variables were the amount of each item (\emph{i-th} co-substrate) to be purchased from each supplier (\emph{j-th} farmer, agro-industrial facility) to maximize the potential profit that may arise from an optimal in-plant management of the purchased quantities. In other words, the aim was to find the potentially optimal co-digestion diet, designed to be constant over the 'purchase horizon' (\emph{N} days of operation). The items are grouped in \emph{R} legislative clusters, subject to \emph{RL$_r$} limits.
Within N, no time domain for the decision variables is considered.
For the simplified approach, the problem was formulated as follows:

\begin{eqnarray*}
\max_{x\in\chi}&&J=\sum_{i=1}^{n}\sum_{j=1}^{m} [p_{ch4} BMP_{_\infty,i} VS_{i} \frac{x_{i,j}}{N} - (C_{Pi,j}+C_{Ti,j})\frac{x_{i,j}}{N})]\\
\mbox{s.t.} && \frac{\sum_{j=1}^{m}\sum_{i=1}^{n\in R} x_{i,j}}{\sum_{j=1}^{m}\sum_{i=1}^{n} x_{i,j}} \lessgtr RL_{r} \ \forall r=1,\dots,R\\
&& x_{i,j} \leq A_{i,j} \ \forall i=1,\dots,n,\ \forall j=1,\dots,m\\
&& k_{min}\leq \frac{\sum_{i=1}^{n}\sum_{j=1}^{m} k_i x_{i,j}}{\sum_{i=1}^{n}\sum_{j=1}^{m} x_{i,j}}\leq k_{max}\ \forall k=1,\dots,K\,\\
&& HRT_{min}\leq \frac{V_RN}{\sum_{i=1}^{n}\sum_{j=1}^{m} x_{i,j}}\leq HRT_{max}\\
&& x_{i,M} = A_{i,M}\ \forall i=1,\dots,n\\
&& x_{i,j}\geq 0\ \forall i=1,\dots,n\ \forall j=1,\dots,m\\
\end{eqnarray*}

where J (\texteuro d$^{-1}$) is the profit, $x_{i,j}$ are the amount of purchased item ($ton_{FM}$ where FM stands for "fresh matter"); $C_{Pi,j}$ and $C_{Ti,j}$ are the specific purchase (\texteuro ton$_{FM_{i}}^{-1}$) and transportation costs (\texteuro ton$_{FM_{i}}^{-1}$) respectively; $V_R$ is the reactor volume ($m^3$), N is the number of days in the 'purchase horizon' (d), $A_{i,j}$ are the ubber bounds for sellers' availabilities; $p_{ch4}$ is the biomethane selling price (\texteuro Nm$_{ch4}^3$) and $VS_i$ (ton$_{VS_{i}}$ton$_{FM_{i}}^{-1}$) are the content of organic matter in the purchased item (the rest being water amd inorganic matter). 
$HRT_{min}, HRT_{max}, k_{min}, k_{max}, RL_w$ are scalar boundaries. $k_i$ are the substrate characteristics subject to bound constraints over the purchased mixture i.e. k = \{Total Kjeldahl Nitrogen (TKN), VS, lipid content, carbon/nitrogen content (C/N)\}.
Unitary density for all items is considered.
The total influent flowrate (m$_{FM}^3d^{-1}$) is thus defined as:
\begin{equation}
    Q_{tot,in} = \frac{\sum_{i=1}^{n}\sum_{j=1}^{m}x_{i,j}}{N}
\end{equation}
Consequently, HRT (d) is defined as: 
\begin{equation}
    HRT = \frac{V_R}{Q_{tot,in}}
\end{equation}
The transportation costs $C_{Ti,j}$ were computed as:
\begin{equation}
    C_{Ti,j} = \frac{FT_i + d_j}{20}
\end{equation}
where $d_j$ is the distance of the \emph{j-th} seller to the plant location and '20' ton$_{FM}$ is the nominal load of transportation to which FT$_i$ refers.
Also, if some stock is already/still present at the plant (player '\emph{M}') at the beginning of the 'purchase horizon', a hard constraint is enforced to guarantee its consumption in the next operating regime.

For the more realistic approach, the problem (later referred as 'LP') was formulated as follows for each HRT value:
\begin{eqnarray*}
\max_{x\in\chi}&&J=\sum_{i=1}^{n}\sum_{j=1}^{m} [p_{ch4} BMP_{i} VS_{i} \frac{x_{i,j}}{N} - (C_{Pi,j}+C_{Ti,j})\frac{x_{i,j}}{N})]\\
\mbox{s.t.} && \frac{\sum_{j=1}^{m}\sum_{i=1}^{n\in R} x_{i,j}}{\sum_{j=1}^{m}\sum_{i=1}^{n} x_{i,j}} \lessgtr RL_{r} \ \forall r=1,\dots,R\\
&& x_{i,j} \leq A_{i,j} \ \forall i=1,\dots,n,\ \forall j=1,\dots,m\\
&& k_{min}\leq \frac{\sum_{i=1}^{n}\sum_{j=1}^{m} k_i x_{i,j}}{\sum_{i=1}^{n}\sum_{j=1}^{m} x_{i,j}}\leq k_{max}\ \forall k=1,\dots,K\,\\
&& \frac{\sum_{i=1}^{n}\sum_{j=1}^{m} (BMP_{\infty,i}-BMP_{i}) x_{i,j}}{\sum_{i=1}^{n}\sum_{j=1}^{m} BMP_{i} x_{i,j}}\leq BMP_{res,max}\\\
&& HRT = \frac{V_RN}{\sum_{i=1}^{n}\sum_{j=1}^{m} x_{i,j}}\\
&& x_{i,M} = A_{i,M}\ \forall i=1,\dots,n\\
&& x_{i,j}\geq 0\ \forall i=1,\dots,n\ \forall j=1,\dots,m\\
\end{eqnarray*}

where each $BMP_i$ is a function of HRT.
